\documentclass[11pt,letterpaper,twoside]{article}
\usepackage[english]{babel}
\usepackage{amssymb,amsmath}
\usepackage{fancyhdr}
\usepackage{graphicx}
\usepackage{wrapfig}

 \oddsidemargin  0in \evensidemargin 0in
 \topmargin   -0.25in \headheight 0.25in \headsep 0.25in
 \textwidth   6.5in \textheight 8.75in \marginparsep 0pt \marginparwidth 0pt
 \parskip 1ex  \parindent 0ex \footskip 20pt

\newfont{\bssten}{cmssbx10}
\newfont{\bssnine}{cmssbx10 scaled 900}

%%%%%%%%%%%%%%%%%%%%%%%
\newcommand{\whatizit}{Lecture 18: Inference for a Difference of Two Proportions}
%%%%%%%%%%%%%%%%%%%%%%%


\pagestyle{fancy}  
\fancyhead{\bssnine STOR 155,  \whatizit}
\fancyhead[RE]{} \fancyhead[LO]{}
\fancyhead[LE]{\bssnine \thepage} \fancyhead[RO]{\bssnine \thepage}
\lfoot{} \cfoot{} \rfoot{}   


\newcommand{\var}{\mathrm{Var}}



\begin{document}



\thispagestyle{empty} \vspace*{-0.75in}

{\bssten STOR 155: Introduction to Data Models and Inference \hfill June 17, 2024 \\
Prof. Will Lassiter  \hfill Page 1 of \pageref{totalpag}}
\vspace{10pt}
\begin{center} {{\Large \bf \whatizit}} \end{center}

Comparing two population proportions: \vspace{400pt}

Example: Do a higher proportion of male college students engage in binge drinking behavior than female students?

\begin{itemize}

\item We randomly sampled $n_1= 5348$ male college students and found that $X_1= 1392$ engage in binge drinking behavior.

\item We randomly sampled $n_2= 8471$ female college students and found that $X_1= 1748$ engage in binge drinking behavior.

\end{itemize}

\newpage

Sampling distribution for a difference of sample proportions: \vspace{200pt}

Back to example: construct a 95\% confidence interval for the difference in binge drinking rates for males and females. \vspace{150pt}

Hypothesis testing for a difference of proportions:

\newpage

Back to example: Can we conclude that male college students engage in binge drinking behavior at a different rate than female college students at a 1\% level of significance? \vspace{250pt}

Example: Growth in student cell phone use in the four years between 2000 and 2004:

\begin{itemize}

\item In 2000, 43\% of a random sample of 1200 students regularly used a cell phone.

\item In 2004, another sample of 1200 students was taken and 89\% regularly used a cell phone.

\end{itemize}

Did cell phone use more than double in 4 years?


\label{totalpag}
\end{document}
